\documentclass[11pt]{article}
\usepackage[utf8]{inputenc}
\usepackage{amsmath, amssymb}
\usepackage{enumitem}
\usepackage{graphicx}
\usepackage{geometry}
\usepackage{hyperref}
\usepackage{xcolor}

\geometry{margin=1in}

\begin{document}

\section{Approximation Algorithms}

Approximation algorithms don't always return an optimal solution, but their solutions are often close to optimal. They always run in polynomial time, and they are often relatively simple.

\subsection{Vertex Cover}

Consider a vertex cover problem. The goal is to return a subset of vertices that covers all edges in a graph. A natural greedy algorithm for vertex cover is to repeatedly select a vertex that covers the most edges and add it to the solution.

However, this greedy algorithm is not always optimal. A better approach is to select an arbitrary edge, add both of its endpoints to the solution, and remove them from the graph. This algorithm has an approximation guarantee of 2.

\subsection{Load Balancing}

Consider a load balancing problem. The goal is to assign jobs to machines to minimize the maximum load. A greedy algorithm for load balancing is to assign each job to the machine with the smallest current load.

This algorithm has an approximation guarantee of 2. We can improve the analysis to show that the algorithm is a $(2 - \frac{1}{m})$-approximation algorithm.

\end{document}